\begin{mistake}{2026-07-15}{04}

\begin{problemcontent}
设
\[
f(x)=\lim_{n\to\infty}\frac{\ln(e^n+x^n)}{n}\quad (x>0),
\]
则
\[
F(x)=\int_1^x f(t)\,dt
\]
在区间 \([1,+\infty)\) 上（ ）。

A. 不连续 \quad
B. 连续但不可导 \quad
C. 可导但二阶不可导 \quad
D. 二阶可导
\end{problemcontent}

\begin{solutioncontent}
先求极限函数。对 \(e^n+x^n\) 中增长较快项分类：
当 \(0<x<e\) 时，
\[
e^n+x^n=e^n\left(1+\left(\frac{x}{e}\right)^n\right),
\]
故
\[
f(x)=\lim_{n\to\infty}\left[1+\frac1n\ln\left(1+\left(\frac{x}{e}\right)^n\right)\right]=1.
\]
当 \(x=e\) 时，
\[
f(e)=\lim_{n\to\infty}\frac{\ln(2e^n)}{n}=1.
\]
当 \(x>e\) 时，
\[
e^n+x^n=x^n\left(1+\left(\frac{e}{x}\right)^n\right),
\]
故
\[
f(x)=\lim_{n\to\infty}\left[\ln x+\frac1n\ln\left(1+\left(\frac{e}{x}\right)^n\right)\right]=\ln x.
\]
所以在 \([1,+\infty)\) 上，
\[
f(x)=
\begin{cases}
1,&1\le x\le e,\\
\ln x,&x>e.
\end{cases}
\]
于是
\[
F(x)=\int_1^x f(t)\,dt
=
\begin{cases}
x-1,&1\le x\le e,\\
\displaystyle \int_1^e1\,dt+\int_e^x\ln t\,dt,&x>e.
\end{cases}
\]
由 \(\int \ln t\,dt=t\ln t-t+C\)，得
\[
F(x)=
\begin{cases}
x-1,&1\le x\le e,\\
x\ln x-x+e-1,&x>e.
\end{cases}
\]
在 \(x=e\) 处，
\[
\lim_{x\to e^-}F(x)=e-1,\qquad
\lim_{x\to e^+}F(x)=e-1,
\]
故 \(F\) 连续。又
\[
F'(x)=
\begin{cases}
1,&1<x<e,\\
\ln x,&x>e,
\end{cases}
\]
且 \(F'_-(e)=1,\ F'_+(e)=\ln e=1\)，故 \(F\) 在 \(e\) 处可导。
但
\[
F''(x)=
\begin{cases}
0,&1<x<e,\\
\frac1x,&x>e,
\end{cases}
\]
在 \(x=e\) 处左右极限分别为 \(0\) 和 \(\frac1e\)，不相等，所以 \(F\) 在 \(e\) 处二阶不可导。

故答案为 \(\boxed{\text{C}}\)。
\end{solutioncontent}

\mistakeknowledge{
\begin{itemize}
  \item \textbf{核心公式/定理}：\(\lim_{n\to\infty}\frac1n\ln(a^n+b^n)=\ln\max\{a,b\}\ (a,b>0)\)；若 \(f\) 连续，则 \(F(x)=\int_a^x f(t)\,dt\) 可导且 \(F'(x)=f(x)\)。
  \item \textbf{易错点}：不能把 \(\ln(e^n+x^n)\) 拆成 \(\ln e^n+\ln x^n\)；判断 \(F\) 二阶可导要看 \(F'=f\) 是否可导。
  \item \textbf{关联知识}：连续的被积函数生成的变上限积分函数一定可导；若 \(f\) 连续但在某点不可导，则对应的 \(F\) 通常可导但二阶不可导。
\end{itemize}}

\end{mistake}
